% ---------------------------------------------------------------------------
% content/links.tex -- the anchor and label conventions.
% ---------------------------------------------------------------------------

\linkedsection{anchors}{Anchors and Labels}

Two different things connect the Hebrew pages to the English ones, and it helps to keep them apart.

\begin{definition}
\label{def:anchor}\hypertarget{target:def:anchor}{}%
An \emph{anchor} is a named position in a PDF, created by \verb|\hypertarget{name}{}| and reached by \verb|\hyperlink{name}{text}|. A \emph{label} is a LaTeX name, created by \verb|\label{name}|, from which \verb|\ref| and \verb|\pageref| print a number and a page.
\end{definition}

Anchors are what a click follows; labels are what gets printed.
The Hebrew document can use both across the document boundary: labels through \texttt{xr-hyper}, and anchors through the merge step, since an anchor's name is stored in the PDF.
The template's convention is to give every object the Hebrew pages may refer to both an anchor and a label, named as in \Cref{tab:naming}.

\begin{table}[htb]
    \centering
    \caption{Naming conventions. The Hebrew document reads labels with the prefix \texttt{en-}.}
    \label{tab:naming}\hypertarget{target:tab:naming}{}%
    \begin{tabular}{@{}>{\raggedright\arraybackslash}p{0.2\textwidth}>{\raggedright\arraybackslash}p{0.4\textwidth}>{\raggedright\arraybackslash}p{0.3\textwidth}@{}}
        \toprule
        Object & Anchor (for \verb|\hyperlink|) & Label (for \verb|\ref*|, \verb|\pageref*|) \\
        \midrule
        Chapter & \texttt{chapter:<key>} & \texttt{chapter:<key>} \\
        Section & \texttt{target:sec:<key>} & \texttt{toc:sec:<key>} \\
        Theorem, figure, table, definition & \texttt{target:thm:<key>}, \texttt{target:fig:<key>}, \ldots & \texttt{thm:<key>}, \texttt{fig:<key>}, \ldots \\
        Bibliography & \texttt{chapter:bib} & \texttt{toc:bibliography} \\
        Bibliography entry & its bib key & \\
        Hebrew summary, abstract & \texttt{chapter:hebrew:summary}, \texttt{chapter:hebrew:abstract} & \texttt{toc:hebrew:summary}, \texttt{toc:hebrew:abstract} \\
        \bottomrule
    \end{tabular}
\end{table}

\linkedsection{sections}{Chapters and Sections}
\label{sec:sections}

Every chapter in \texttt{main.tex} follows the same block.
The \verb|\phantomsection| gives hyperref a correct bookmark, the \verb|\hypertarget| is the anchor the Hebrew TOC links to, and the \verb|\label| is where it reads the chapter number and page.

\begin{lstlisting}
\cleardoublepage
\phantomsection
\hypertarget{chapter:build}{}%
\chapter{The Build Process}
\label{chapter:build}
% ---------------------------------------------------------------------------
% content/build.tex -- how compile_thesis.sh works and why.
% ---------------------------------------------------------------------------

\linkedsection{script}{What the Script Does}

\texttt{compile\_thesis.sh} needs TeX Live (or MacTeX) with \texttt{latexmk}, \texttt{pdflatex}, \texttt{xelatex} and \texttt{biber}, and Python~3 with the \texttt{pikepdf} package.
It runs the three steps of \Cref{fig:pipeline}, all inside the git-ignored folder \texttt{build/}.

\begin{figure}[htb]
\hypertarget{target:fig:pipeline}{}%
    \centering
    \resizebox{\textwidth}{!}{%
    \begin{tikzpicture}[
        box/.style={draw,rounded corners,minimum height=9mm,minimum width=27mm,align=center,font=\small},
        arr/.style={-{Stealth[length=2.5mm]},thick}]
        \node[box] (en) {\texttt{main.tex}\\pdfLaTeX + biber};
        \node[box,below=9mm of en] (he) {\texttt{hebrew\_pages.tex}\\XeLaTeX + biber};
        \node[box,right=22mm of en] (enpdf) {\texttt{build/main.pdf}};
        \node[box,right=22mm of he] (hepdf) {\texttt{build/hebrew\_pages.pdf}};
        \node[box,right=22mm of $(enpdf)!0.5!(hepdf)$] (merge) {\texttt{merge\_thesis.py}};
        \node[box,right=12mm of merge] (out) {\texttt{build/thesis.pdf}};
        \draw[arr] (en) -- (enpdf);
        \draw[arr] (he) -- (hepdf);
        \draw[arr] (enpdf) -- (merge);
        \draw[arr] (hepdf) -- (merge);
        \draw[arr] (merge) -- (out);
        \draw[arr,dashed] (en.south) to[bend right=25] node[left,font=\scriptsize,align=right] {page numbers,\\\texttt{LastPage}} (he.north);
        \draw[arr,dashed] (he.north east) to[bend right=25] node[right,font=\scriptsize,align=left] {Hebrew\\page numbers} (en.south east);
    \end{tikzpicture}}
    \caption{The build. Solid arrows are the three steps of the script; dashed arrows are the labels each document reads from the other's auxiliary file, which is why the first two steps are repeated until both settle.}
    \label{fig:pipeline}
\end{figure}

\begin{enumerate}
    \item \texttt{latexmk -pdf main.tex} produces \texttt{build/main.pdf}. \texttt{latexmk} runs biber and repeats pdfLaTeX until the references are stable.
    \item \texttt{latexmk -xelatex hebrew\_pages.tex} produces \texttt{build/hebrew\_pages.pdf}.
    \item \texttt{merge\_thesis.py} appends the Hebrew pages to the English PDF and rewires the cross-document links (\Cref{sec:merge}). The result is \texttt{build/thesis.pdf}.
\end{enumerate}

\texttt{./compile\_thesis.sh clean} deletes \texttt{build/}. A clean build takes a few minutes; later builds only redo what changed.

\linkedsection{settle}{Why the First Two Steps Repeat}

Each document reads labels from the other.
The Hebrew document continues the page numbering after the English \texttt{LastPage} label and prints the English chapter numbers and pages in its table of contents.
The English document, in turn, needs the page numbers of the Hebrew summary and abstract for its own table of contents.
Both documents load the \texttt{xr-hyper} package, which imports the labels of another document's \texttt{.aux} file under a prefix: \texttt{en-} for English labels read by the Hebrew document, \texttt{he-} for Hebrew labels read by the English one.

Because each side depends on the other, the script compiles English, then Hebrew, then checks whether either \texttt{.aux} file changed, and repeats the pair until neither does.
On a clean build the first English pass prints page~0 for the Hebrew entries, since no Hebrew \texttt{.aux} exists yet; the second pass fixes it.
From the second pass on, the script passes \texttt{-g} to \texttt{latexmk}, because \texttt{latexmk} does not know that the other document's \texttt{.aux} file is a dependency and would otherwise report that there is nothing to do.

\linkedsection{troubleshooting}{Fonts, Overleaf and Troubleshooting}

\paragraph{Fonts.} The Hebrew pages use \emph{Linux Libertine O} for Latin text, \emph{David CLM} for Hebrew and \emph{Latin Modern Roman} for Latin fragments inside Hebrew paragraphs. All three ship with TeX Live (packages \texttt{libertine}, \texttt{culmus}, \texttt{lm}), but XeTeX on macOS only finds a font by name if macOS itself knows it. If the build stops with \texttt{The font "Linux Libertine O" cannot be found}, run \texttt{./compile\_thesis.sh install-fonts} once; it copies the font files into \texttt{\textasciitilde/Library/Fonts}. On Linux and on Overleaf this is not needed.

\paragraph{Overleaf.} Overleaf cannot run the merge script. Compile \texttt{main.tex} with pdfLaTeX and \texttt{hebrew\_pages.tex} with XeLaTeX (two projects, or one project with the main file switched), download both PDFs and merge locally with \texttt{python3 merge\_thesis.py english.pdf hebrew.pdf out.pdf}.

\paragraph{Symptoms and causes.}
\begin{itemize}
    \item \emph{Hebrew Summary~0 in the English TOC}: the English part was not rebuilt after the Hebrew part. Run the script again rather than \texttt{latexmk} by hand.
    \item \emph{Merge warning about names that exist in neither PDF}: a Hebrew \verb|\hyperlink| has no matching anchor in the English source. See \Cref{sec:merge}.
    \item \emph{Nonsense errors such as \texttt{Missing \textbackslash begin\{document\}} in \texttt{build/main.aux}}: a build was interrupted or two ran at once and an \texttt{.aux} file is truncated. Run \texttt{./compile\_thesis.sh clean} and build again.
    \item \emph{biber: extracting arm64 binary with lipo failed} (macOS, after a system update): the universal \texttt{biber} shipped with MacTeX unpacks itself with \texttt{lipo}, which newer command line tools reject. Extract the slice once by hand, \texttt{lipo -thin arm64 \$(which biber) -output \textasciitilde/bin/biber}, and put \texttt{\textasciitilde/bin} first on your \texttt{PATH}.
    \item \emph{Package bidi error about xcolor}: in \texttt{hebrew\_pages.tex}, any package that loads \texttt{xcolor} must come before \texttt{bidi} and \texttt{polyglossia}. The template loads \texttt{xcolor} first for this reason.
\end{itemize}

\end{lstlisting}

Sections use the helper \verb|\linkedsection|, defined in \texttt{main.tex}:

\begin{lstlisting}
\linkedsection{script}{What the Script Does}
% expands to
\section[What the Script Does]{\hypertarget{target:sec:script}{}What the Script Does}
\label{toc:sec:script}
\end{lstlisting}

The anchor is placed \emph{inside} the heading on purpose. An anchor written on the line before \verb|\section| is set at the end of the preceding paragraph, so whenever the section happens to start a new page, a click lands at the bottom of the previous one.
Add your own \verb|\label{sec:...}| after \verb|\linkedsection| for \verb|\Cref|, as this section does.

An unnumbered chapter, such as the conclusion of this document, has no number to read, so it needs only the anchor and a label, plus \verb|\addcontentsline| to appear in the English TOC.
\Cref{alg:chapter} summarises the steps for adding a chapter that both tables of contents can see.

\begin{algorithm}[htb]
\caption{Adding a chapter}
\label{alg:chapter}
\KwIn{a chapter title, a short key, a Hebrew title}
write the chapter in \texttt{content/<key>.tex}, using \texttt{\textbackslash linkedsection} for its sections\;
copy the chapter block into \texttt{main.tex} with \texttt{chapter:<key>} as anchor and label\;
\ForEach{section in the chapter}{
    add a line to the Hebrew TOC with \texttt{toc:sec:<key>} and \texttt{target:sec:<key>} (\Cref{sec:hebrewtoc})\;
}
add a line for the chapter itself with \texttt{chapter:<key>}\;
run \texttt{./compile\_thesis.sh} and check that the merge step prints no warning\;
\end{algorithm}

\linkedsection{objects}{Theorems, Figures, Tables and Definitions}
\label{sec:objects}

Anything the Hebrew text should be able to point at gets an anchor right after its label.
For a theorem, the anchor goes at the start of the body; for a figure or table, right after \verb|\begin{figure}| or next to the \verb|\label|.

\begin{lstlisting}
\begin{restatable}{theorem}{MainTheorem}
\label{thm:main}\hypertarget{target:thm:main}{}%
Every object in the English part that carries ...
\end{restatable}

\begin{figure}[htb]
\hypertarget{target:fig:pipeline}{}%
    ...
    \label{fig:pipeline}
\end{figure}
\end{lstlisting}

The theorem above uses \texttt{restatable} from the \texttt{thmtools} package: it is stated in the introduction and restated here with \verb|\MainTheorem*|, keeping its number, so the two statements cannot drift apart.

\MainTheorem*

\begin{proof}
Let $x$ be an object with anchor \texttt{target:x}.
The Hebrew source writes \verb|\hyperlink{target:x}{...}|, which XeLaTeX turns into a link whose destination is the \emph{name} \texttt{target:x}.
Inside the Hebrew PDF alone this name is undefined.
After the merge, the combined PDF contains both the link and, from the English half, the anchor of that name (\Cref{sec:merge}).
The merge script then replaces the name in the link by the page and position of the anchor, so the link no longer depends on the name table, which some viewers discard when they re-save a file.
\end{proof}

\linkedsection{merge}{The Merge Step}
\label{sec:merge}

\texttt{merge\_thesis.py} appends the pages of \texttt{build/hebrew\_pages.pdf} to \texttt{build/main.pdf}, collects the named anchors of both files, and rewrites every link that refers to a name into an explicit destination, as in the proof above.
Names defined in both files are taken from the English half, except the \texttt{chapter:hebrew:*} anchors, which are taken from the Hebrew half.

The script prints how many links it resolved and, if any Hebrew link uses a name that neither document defines, a warning listing those names.
A clean thesis prints no warning.
The fix is to add the missing \verb|\hypertarget| in the English source.
As a last resort, the \texttt{ALIASES} table at the top of the script can map a name to one of hyperref's automatic anchors, such as \texttt{section.2.1} or \texttt{theorem.3.1}; those names change whenever numbering changes, so prefer a real anchor.

Citations need no anchors of their own: both documents load the same \texttt{.bib} file with the same \texttt{biblatex} options, and after the merge a Hebrew citation link resolves to the entry in the English bibliography (\Cref{sec:hcite}).
In addition, \texttt{main.tex} gives every bibliography entry an anchor named by its bib key, so \verb|\hyperlink{Karp72}{...}| works from either side.
